% The instrument in series with the board's own supply, plus three marker pins.
\begin{circuitikz}[font=\sffamily\scriptsize, line width=0.6pt]
% ---------------- the board's own 3.3 V rail
\draw (0,-1.0) to[battery1] (0,0);
\draw (0,-1.0) -- (0,-1.6) node[ground]{};
\node[anchor=east, align=right] at (-0.55,-0.5) {3V3 from the\\board regulator};
\draw (0,0) -- (1.2,0);

% ---------------- the measurement jumper, opened
\draw (1.2,0) to[short, -o] (1.7,0);
\draw (3.1,0) to[short, o-] (5.5,0);
\node[anchor=south, align=center] at (2.4,0.25) {measurement jumper\\opened, item 6};

% ---------------- the instrument, bridging the opened jumper
\draw (1.7,0) -- (1.7,-1.6) to[ammeter] (3.1,-1.6) -- (3.1,0);
\draw[dashed, draw=linecol] (1.25,-2.25) rectangle (3.55,-1.05);
\node[anchor=north, align=center] at (2.4,-2.35) {nRF-PPK2\\current channel};

% ---------------- the part under measurement
\node[draw=linecol, fill=hwcol, minimum width=30mm, minimum height=14mm,
      align=center] (mcu) at (7.0,0) {STM32H7A3ZI\\{\tiny VDD, 1 A maximum}};
\draw (5.7,-0.7) -- (5.7,-1.3) node[ground]{};

% ---------------- three marker pins into the digital inputs
\draw[dashed, draw=linecol] (5.6,-3.25) rectangle (8.6,-2.2);
\node[anchor=north west, font=\sffamily\tiny] at (5.62,-3.28) {nRF-PPK2 digital inputs};
\foreach \x/\name in {6.8/D0, 7.4/D1, 8.0/D2}{
  \node[draw=black!70, fill=gray!70, minimum width=2.6mm, minimum height=2.6mm]
        (t\name) at (\x,-2.6) {};
  \node[font=\tiny, anchor=north] at (\x,-2.78) {\name};
  \draw[cable, draw=orange!70!black] (\x,-0.7) -- (\x,-2.47);}
\node[draw=black!70, fill=gray!70, minimum width=2.6mm, minimum height=2.6mm]
      (tg) at (6.0,-2.6) {};
\node[font=\tiny, anchor=north] at (6.0,-2.78) {GND};
\draw[cable, draw=black!55] (6.0,-0.7) -- (6.0,-2.47);

\node[umlnote, text width=54mm, anchor=north west] at (-1.8,-3.9)
  {One instrument, two kinds of channel, one timebase. The board drives the digital inputs and they never drive the board. The probe is still powered from USB and is not in the measured path.};
\end{circuitikz}
